\section{Set covering with precedence constraints} \label{sec:SC}

% Similarly as in previous section we assume our input instance to be unweighted. Note that the analysis of the following procedures works also in the weighted setting, where every element $u\in\cU$ has a weight $w(u)$ and the goal is to maximize the total weight of covered elements or cover at least a given total weight of elements with minimal size of the cover. Since this generalization is straightforward we omit it for clarity.

Having obtained appropriate approximation algorithms for $\ProblemMDPCS$ and $\ProblemBMDPCS$, we now turn our attention towards approximating various set covering problems. In particular, we show how to obtain a bicriteria $\br{\bigo\br{\sqrt{m\cdot \log n}}, 1}$-approximation for $\ProblemBPCSC$. This can be then turned into a $\br{\bigo\br{\sqrt{m\cdot \log n}}, 1}$-approximation for $\ProblemPCSC$ and a $\br{\bigo\br{\sqrt{m\cdot \log n}}, O\br{1}}$-approximation for $\ProblemPCMSSC$. At the end of the section we also show several better guarantees for inforest and outforest precedence constraints.
\subsection{$\ProblemBPCSC$ via $\ProblemBMDPCS$} \label{subsection:BPCSC}

We start with the budgeted variant, which will become a useful primitive for the fractional variants.
The approximation algorithm is the greedy procedure shown in Algorithm~\ref{alg:BPCSC}.
For the analysis, let $\cC_i$ and $\cA_i$, $i\in[l]$, be the sets $\cC$ and $\cA$ obtained in the $i$-th iteration; $\cC_0=\emptyset$.
Note that $\cov\br{\cA_i,\cU\setminus\cov\br{\cC_{i-1}}}$ is the set of element in $\cU$ that are being covered in the $i$-th iteration. %, and for each $u\in\cov\br{\cA_i,\cU\setminus\cC_{i-1}}$ denote $c(u):=\Delta\br{\cA_i,\cU\setminus\cC_{i-1}}$.
\begin{algorithm}[htb!]
\DontPrintSemicolon
\LinesNumbered
\caption{An $\br{\alpha_{m,n},1}$-approximation algorithm for $\ProblemBPCSC$}
\label{alg:BPCSC}
\SetKwProg{Proc}{procedure}{}{}
\Proc{$\ProcedureBPCSC(\cU, \cS, \preceq, \budget)$}{
  $\cR\gets\brc{r_S:S\in\cS}$\;
  $\cB\gets\brc{b_u:u\in\cU}$\;
  $\cA\gets\emptyset$\;

  \ForEach{$S\in\cS$}{
    $D_S\gets\closure{S}$\;
    $A_S\gets\brc{r_T:T\in D_S}\cup\brc{b_u:u\in\cov\br{D_S}}$\;
    $\cA\gets\cA\cup\brc{A_S}$\;
  }

  $\cC_{\mathrm{best}}\gets\emptyset$\;
  \For{$t\gets 1$ \KwTo $\spr{\cov\br{\cS}}$}{
    $\cA_t\gets$ Run the $\alpha_{m,n}$-approximation algorithm for $\ProblemPRBSC$ on $(\cR,\cB,\cA,t)$\;
    $\cC_t\gets\bigcup_{A_S\in\cA_t}D_S$\;
    \If{$\spr{\cC_t}\leq\alpha_{m,n}\cdot\budget \land \spr{\cov\br{\cC_t}}>\spr{\cov\br{\cC_{\mathrm{best}}}}$}{
      $\cC_{\mathrm{best}}\gets\cC_t$\;
    }
  }
  \Return $\cC_{\mathrm{best}}$\;
}
\end{algorithm}


First we bound the size of the set returned by the algorithm (cf. Lemma~\ref{lem:BPCSC-cost}) and then we estimate how much it covers (cf. Lemma~\ref{lem:BPCSC-coverage}).
\begin{lemma}\label{lem:BPCSC-cost}
    Let $\cC$ be the cover returned by Algorithm~\ref{alg:BPCSC}. Then, $\spr{\cC} \leq \br{\sqrt{m\cdot H_n}+1}\cdot \budget$.
\end{lemma}
\begin{proof}
    We consider two cases:
    \begin{enumerate}
        \item If $\budget\geq \sqrt{m/H_n}$, then for $\cC=\cS$ it holds: $\spr{\cC} = m \leq \sqrt{m\cdot H_n}\cdot \budget$.
        \item Else, when $\budget < \sqrt{m/H_n}$, we have that before the last iteration, that is in iteration $l-1$, of the while loop, $\spr{\cC_{l-1}} < \sqrt{m\cdot H_n}\cdot \budget$. Since in the last iteration we add to $\cC_{l-1}$ sets in a collection $\cA_l$ such that $\spr{\cA_l} \leq \budget$, we have that 
        \[
        \spr{\cC_l} \leq \sqrt{m\cdot H_n}\cdot \budget + \budget = \br{\sqrt{m\cdot H_n}+1}\cdot \budget.
        \]
    \end{enumerate}
\end{proof}

\begin{lemma}\label{lem:BPCSC-coverage}
    Let $\cC$ be the cover returned by Algorithm~\ref{alg:BPCSC}. Then, $\spr{\cov\br{\cC}} \geq \spr{\cov\br{\cCopt}}$.
\end{lemma}
\begin{proof}
    If $\budget \geq \sqrt{m/H_n}$, then the algorithm returns $\cC_l=\cS$ and trivially covers at least as many elements as $\cCopt$.

    Otherwise, let $q=\spr{\cov\br{\cCopt}}$. If $q=0$, the claimed coverage guarantee is immediate, so assume $q>0$. We analyze the loop up to the first iteration $t$ for which $\spr{\cov\br{\cC_t}}\geq q$ and show that this iteration occurs before the algorithm reaches its cost threshold.

    Let $\cC_i$ and $\cA_i$ be the cover and the family added in iteration $i$, with $\cC_0=\emptyset$, and define $R_i=\cU\setminus\cov\br{\cC_i}$ for all $i\geq 0$. In particular, $R_0=\cU$. For each covered element $u$, let $i(u)$ be its first covering iteration and define its price by
    \[
    c(u)=\frac{\spr{\cA_{i(u)}}}{\spr{\cov\br{\cA_{i(u)},R_{i(u)-1}}}}.
    \]

    Set $k=\spr{\cov\br{\cC_{t-1}}}<q$, and order the $k$ elements of $\cov\br{\cC_{t-1}}$ as $v_1,\ldots,v_k$ by their first covering iteration, breaking ties arbitrarily.
    \begin{claim}\label{lem:cost-per-element}
        For every $j\in[k]$, $c(v_j)\leq\frac{\budget^2}{q-j+1}$.
    \end{claim}
    \begin{proof}
        Let $i=i(v_j)$. At most $j-1$ elements have been covered before iteration $i$, so at least $q-j+1$ elements of $\cov\br{\cCopt}$ remain in $R_{i-1}$. The family $\mathcal{D}_i=\cCopt\setminus\cC_{i-1}$ is precedence-closed in the residual family $\cS\setminus\cC_{i-1}$, has size at most $\spr{\cCopt}\leq\budget$, and covers all elements of $\cov\br{\cCopt}\cap R_{i-1}$. Thus, $\spr{\cov\br{\mathcal{D}_i,R_{i-1}}}\geq q-j+1$. By the $\budget$-approximation guarantee for $\ProblemBMDPCS$,
        \[
        \Delta\br{\cA_i,R_{i-1}}\geq\frac{1}{\budget}\Delta\br{\mathcal{D}_i,R_{i-1}}\geq\frac{q-j+1}{\spr{\mathcal{D}_i}\cdot\budget}\geq\frac{q-j+1}{\budget^2}.
        \]
        Since $c(v_j)=1/\Delta\br{\cA_i,R_{i-1}}$, the claim follows.
    \end{proof}

    By the price definition, summing over exactly the elements covered before iteration $t$ gives
    \[
    \spr{\cC_{t-1}}=\sum_{j=1}^{k}c(v_j)
    \leq\budget^2\sum_{j=1}^{k}\frac{1}{q-j+1}
    =\budget^2\br{H_q-H_{q-k}}
    \leq\budget^2 H_n.
    \]
    Since $\spr{\cA_t}\leq\budget$, it follows that
    \[
    \spr{\cC_t}\leq\budget^2 H_n+\budget
    \leq\br{\sqrt{m\cdot H_n}+1}\cdot\budget.
    \]
    Moreover, $\budget<\sqrt{m/H_n}$ in this case, so $\spr{\cC_{t-1}}\leq\budget^2 H_n<\sqrt{m\cdot H_n}\,\budget$. Hence iteration $t$ is reached before the algorithm's stopping threshold, and the algorithm covers at least $q$ elements. This proves the lemma.
\end{proof}

Lemmas \ref{lem:BPCSC-cost} and \ref{lem:BPCSC-coverage} prove the following.
\begin{theorem} \label{thm:BPCSC}
    There exists a polynomial-time $\br{\sqrt{m\cdot H_n}+1, 1}$-bicriteria approximation algorithm for $\ProblemBPCSC$, where $n=\spr{\cU}$ and $m=\spr{\cS}$. That is, the algorithm returns a solution $\cC$ such that $\spr{\cov\br{\cC}}\geq \OPT\br{\cI}$ and $\spr{\cC}\leq \br{\sqrt{m\cdot H_n}+1}\cdot \budget$.
\end{theorem}

\subsection{$\ProblemPCSC$ via $\ProblemMDPCS$}\label{subsection:PCSC}
Let $k=f\cdot n$.
In order to prove the main theorem of this section consider a greedy procedure shown as Algorithm~\ref{alg:PCSC}.
We note a subtle fact of finishing the loop once $k/2$ elements are covered and not iterating till reaching $k$.
This is because the last set $\cA$ may be of unbounded size but we argue that the size of an optimal solution to the instance from the last iteration is lower-bounded by $\Omega(k)$, which is enough to obtain the required approximation ratio.
\begin{algorithm}[h]
\DontPrintSemicolon
\LinesNumbered
\caption{An $\alpha_{m,n}$-approximation algorithm for $\ProblemPCSC$}
\label{alg:PCSC}
\SetKwProg{Proc}{procedure}{}{}
\Proc{$\ProcedurePCSC(\cU, \cS, \preceq, k)$}{
  $\cR\gets\brc{r_S:S\in\cS}$\;
  $\cB\gets\brc{b_u:u\in\cU}$\;
  $\cA\gets\emptyset$\;

  \ForEach{$S\in\cS$}{
    $D_S\gets\closure{S}$\;
    $A_S\gets\brc{r_T:T\in D_S}\cup\brc{b_u:u\in\cov\br{D_S}}$\;
    $\cA\gets\cA\cup\brc{A_S}$\;
  }

  $\cA'\gets$ Run the $\alpha_{m,n}$-approximation algorithm for $\ProblemPRBSC$ on $(\cR,\cB,\cA,k)$\;
  $\cC\gets\bigcup_{A_S\in\cA'}D_S$\;
  \Return $\cC$\;
}
\end{algorithm}

\begin{theorem} \label{thm:MDPCStoPCSC}
    If there exists a polynomial-time $\gamma$-approximation algorithm for $\ProblemMDPCS$, then there exists a polynomial-time bicriteria $\br{4\gamma/f, 2}$-approximation algorithm for $\ProblemPCSC$.
\end{theorem}
\begin{proof}
Denote by $\cA_i$ the set $\cA$ from the $i$th iteration of Algorithm~\ref{alg:PCSC}, $i\in[l]$.
For brevity let $\cC_i$, $i\in[l]$, be the set $\cC$ obtained in the $i$th iteration, $\cC_i=\bigcup_{j=1}^i \cA_j$.
Let $R_i=\cU\setminus\cov\br{\cC_i}$ for $i\in[l]$ and $R_0=\cU$, $\cC_0=\emptyset$.
Denote by $\cI=(\cU,\cS,\preceq,k)$ any input instance to $\ProblemPCSC$, and let $\cCopt$ be an optimal solution to $\ProblemPCSC$ on $\cI$, i.e., $\cov\br{\cCopt}\geq k$ and $\spr{\cCopt}=\optPCSC{\cI}$.

Note that for each $i\in[l]$, $\cCopt\setminus\cC_{i-1}\neq\emptyset$ because otherwise $\cC_{i-1}\subseteq\cCopt$ which means that $\spr{\cC_{i-1}}\geq\spr{\cCopt}\leq k$ contradicting the fact that the algorithm conducted the $l$-th iteration.
Hence, $\cCopt\setminus\cC_{i-1}$ is a precedence-closed family of density not greater than the density of an optimal solution to $\ProblemMDPCS$ for the input provided in the $i$-th iteration, $i\in[l]$.
Thus, 
\[
\Delta(\cA_i,R_{i-1}) \geq \frac{\Delta(\cCopt\setminus\cC_{i-1},R_{i-1})}{\gamma},
\]
which gives by definition of $\Delta$:
\[
\frac{\spr{\cov\br{\cA_i,R_{i-1}}}}{\spr{\cA_i}} \geq \frac{1}{\gamma}\cdot\frac{\spr{\cov\br{\cCopt,R_{i-1}}}}{\spr{\cCopt\setminus\cC_{i-1}}}, \quad i\in[l],
\]
which further implies:
\begin{equation} \label{eq:Ai}
\spr{\cA_i} \leq \gamma \cdot \frac{ \spr{\cov\br{\cA_i,R_{i-1}}} \cdot \spr{\cCopt\setminus\cC_{i-1}} }{ \spr{\cov\br{\cCopt,R_{i-1}}} }
            \leq \gamma \cdot \frac{ \spr{\cov\br{\cA_i,R_{i-1}}} }{ \spr{\cov\br{\cCopt,R_{i-1}}} } \cdot \spr{\cCopt}.
\end{equation}
For each $i\in[l-1]$ it holds $\spr{\cov\br{\cCopt,R_{i-1}}}\geq k/2$.
Thus,
\[
\sum_{i=1}^{l-1}\spr{\cA_i} \leq \frac{2\gamma}{k} \cdot \spr{\cCopt} \cdot \sum_{i=1}^{l-1}\spr{\cov\br{\cA_i,R_{i-1}}} \leq \frac{2\gamma n}{k} \cdot \spr{\cCopt}.
\]
For the last iteration $\spr{\cov\br{\cA_l,R_{l-1}}}\leq n$ (potentially $\cA_l$ may be of size unbounded by a function of $k$) and $\spr{\cov\br{\cCopt,R_{l-1}}}\geq k/2$.
Hence by \eqref{eq:Ai} we get $\spr{\cA_l}\leq \frac{2\gamma n}{k}\spr{\cCopt}$.
This gives 
\[
\spr{\cC_l}=\sum_{i=1}^l\spr{\cA_i}\leq\frac{4\gamma n}{k}\cdot\spr{\cCopt} = \frac{4\gamma}{f}\cdot\spr{\cCopt}.
\]
Moreover, $\spr{\cov\br{\cC_l}}\geq k/2$, which means the algorithm covers at least half of the required elements.
This proves theorem.
\end{proof}

Additionally, observe that if one has an $\br{\alpha, \beta}$-approximation algorithm for $\ProblemBPCSC$, then one can obtain an $\br{\alpha, \beta}$-approximation algorithm for $\ProblemPCSC$ by simply guessing the optimal budget $\budget^*=\spr{\cCopt}$ and running the $\ProblemBPCSC$ algorithm with budget $\budget^*$. We have the following corollary.
\begin{corollary}\label{cor:BPCSC-to-PCSC}
    If there exists a polynomial-time $(\alpha,\beta)$-approximation algorithm for $\ProblemBPCSC$, then there exists a polynomial-time $(\alpha,\beta)$-approximation algorithm for $\ProblemPCSC$.
\end{corollary}


\subsection{$\ProblemPCMSSC$ via $\ProblemBPCSC$} \label{subsection:PCMSSC}
We show how to use an $(\alpha,\beta)$-approximation algorithm for $\ProblemBPCSC$ to obtain an $\br{\bigo\br{\alpha\cdot\beta^3}, \beta}$-approximation for $\ProblemPCMSSC$. This result seems counterintuitive at first glance, since we are using an algorithm for a problem where the budget denotes the maximal allowed size of the cover to solve a problem where the budget denotes the sum of covering times. Our algorithm and its analysis are a generalization of the approach of \cite{ApproxAlgsForOptDTsAndAdapTSPProblems}. Furthermore, their result is based on techniques used for the \textsc{Minimum Latency Travelling Salesman Problem} \cite{PathsTreesMinimumLatencyTours,KTravelingRepairmenProblem}.

Let $\alpha,\beta\geq 1$ and $a=\frac{3\beta-1}{3\beta-2}$, so $1<a\leq 2$.
Consider the Algorithm~\ref{alg:fPCMSSC}.
% Denote by $\cCopt$ any optimal solution to $\ProblemPCMSSC$ on the input $(\cU, \cS, \preceq, f)$.
Let $l^*=\max\brc{1,\cl{\log_a\spr{\cCopt}}}$, so $a^{l^*-1}\leq\spr{\cCopt}\leq a^{l^*}$.
\begin{algorithm}[h]
\DontPrintSemicolon
\LinesNumbered
\caption{The algorithm for $\ProblemPCMSSC$ based on $\br{\alpha, \beta}$-approximation for $\ProblemBPCSC$}
\label{alg:fPCMSSC}
\SetKwProg{Proc}{procedure}{}{}
Fix a topological ordering of $(\cS,\preceq)$; $\operatorname{ord}(A)$ denotes its restriction to $A\subseteq\cS$.
\Proc{$\ProcedureFPCMSSC(\cU, \cS, \preceq, f)$}{
    $\cQ\gets \emptyset$.

    \For{$1\leq l\leq \max\brc{1,\cl{\log_a{\spr{\cS}}}}$}{
        $U\gets \cU$.

        $\pi_l\gets \emptyset$.
        $P_l\gets \emptyset$.

        \For{$1\leq i\leq l$}{
            $\tau_l^i\gets $ Run the $\br{\alpha, \beta}$-approx. algorithm for $\ProblemBPCSC$ on $(U, \preceq, \cS, \budget=a^{i+1})$.

            $U \gets U \setminus \cov(\tau_l^i)$.

            $\pi_l \gets \pi_l \circ \operatorname{ord}(\tau_l^i\setminus P_l)$.
            $P_l \gets P_l\cup\tau_l^i$.
        }
        $T_l\gets\min\brc{m,\cl{\alpha\cdot\beta\cdot a^l}}$.
        \While{$\spr{\pi_l}<T_l$}{
            $S_l\gets$ the next set in the fixed topological ordering of $\cS$ that is not in $P_l$.
            $\pi_l\gets\pi_l\circ\angl{S_l}$.
            $P_l\gets P_l\cup\brc{S_l}$.
        }

        $R_l\gets\cU\setminus\cov(P_l)$.
        $\sigma_l \gets $ Run the $\br{\alpha, \beta}$-approx. algorithm for $\ProblemBPCSC$ on $(R_l, \preceq, \cS, \budget=a^l)$.
        $\pi_l\gets\pi_l\circ\operatorname{ord}(\sigma_l\setminus P_l)$.

        \If{$\spr{\cov\br{\pi_l,\cU}}\geq f\cdot\spr{\cU}/\beta$}{
            $\cQ \gets \cQ \cup \brc{\pi_l}$.
        }
    }
    \Return $\argmin_{\pi\in \cQ} \brc{\coverageTime\br{\pi,\cU}}$.
}
\end{algorithm}


Each family returned by $\ProblemBPCSC$ is precedence-closed, so the union of the sets already appended is also precedence-closed. When a family is appended, only its new sets are used and they are ordered by $\operatorname{ord}$; every predecessor of a new set is therefore either already in the sequence or earlier in the appended block. The same holds for the padding and the final family, so every sequence considered by the algorithm is precedence-closed and contains no repeated sets.
We start with a lemma that relates the coverage times and the lengths of $\pi_{l^*}$ and $\cCopt$.
    \begin{lemma}\label{lem:pi-spread-cost}
        We have that $\spr{\pi_{l^*}} = \bigo\br{\alpha \cdot \beta}\cdot \spr{\cCopt}$, and if $\spr{\pi_{l^*}}<m$, then $\spr{\pi_{l^*}}=\Omega\br{\alpha\cdot\beta}\cdot\spr{\cCopt}$. Moreover, $\coverageTime\br{\pi_{l^*},\cU} = \bigo\br{\alpha\cdot\beta^3}\cdot \coverageTime\br{\cCopt,\cU}$.
    \end{lemma}
    \begin{proof}
        We observe that since we used an $\br{\alpha, \beta}$-approximation algorithm for $\ProblemBPCSC$ in each iteration of the inner for loop, we have that for each $1\leq i\leq l^*$, $\spr{\tau_{l^*}^i} \leq \alpha \cdot a^{i+1}$. Thus, the length after padding satisfies:
        \begin{equation} \label{eq:pi}
        \spr{\pi_{l^*}} \leq \max\brc{\alpha\cdot \sum_{i=1}^{l^*} a^{i+1},\min\brc{m,\cl{\alpha\cdot\beta\cdot a^{l^*}}}}\leq \frac{\alpha\cdot a^{l^*+3}}{a-1}\leq \frac{\alpha\cdot a^{4}}{a-1}\cdot\spr{\cCopt}=\bigo\br{\alpha\cdot\beta}\cdot\spr{\cCopt}.
        \end{equation}
        Here we used $\cl{x}\leq x+1$, $\alpha a^{l^*}\geq1$, $\beta\leq1/(a-1)$, and $a^{l^*}\leq a\spr{\cCopt}$. If $\spr{\pi_{l^*}}<m$, the padding target is not capped by $m$, and hence $\spr{\pi_{l^*}}\geq\alpha\cdot\beta\cdot a^{l^*}\geq\alpha\cdot\beta\cdot\spr{\cCopt}$. If $\spr{\pi_{l^*}}=m$, then the prefix contains every set in $\cS$; this case will be handled after the latency bound.

        For every $i\in [l^*]$, denote $n_i^*=\spr{\cov\br{\pref{\cCopt}{a^{i}}}}$ and note that $n_{l^*}^* = \spr{\cov\br{\cCopt}}$. Accordingly, for every $i\in [l^*]$ let $n_i=\spr{\cov\br{\angl{\tau_{l^*}^1,\dots, \tau_{l^*}^i}}}$. Additionally, define $n_0=n_0^*=0$. We have:
        \begin{align*}
            \coverageTime\br{\pi,\cU}&\leq \sum_{i=1}^{l^*}\br{n_i-n_{i-1}}\cdot \sum_{j=1}^{i}\alpha\cdot a^{j+1} + \br{n-n_{l^*}}\cdot \spr{\pi_{l^*}}
            \\&\leq
            \sum_{i=1}^{l^*}\br{n_i-n_{i-1}}\cdot \frac{\alpha\cdot a^{i+2}}{a-1} + \br{n-n_{l^*}}\cdot \spr{\pi_{l^*}}
            \\&=
            \sum_{i=1}^{l^*}\br{\br{n-n_{i-1}}-\br{n-n_i}}\cdot \frac{\alpha\cdot a^{i+2}}{a-1} + \br{n-n_{l^*}}\cdot \spr{\pi_{l^*}}
            \\&\leq
            \sum_{i=0}^{l^*}\br{n-n_i}\cdot \frac{\alpha\cdot a^{i+3}}{a-1} =: Q.
        \end{align*}
        where the last inequality is due to \eqref{eq:pi}.
        Moreover, we have that:
        \begin{align*}
        2\cdot\coverageTime\br{\cCopt,\cU}&\geq a^{l^*-1}\br{n-n_{l^*-1}^*}+ \sum_{i=1}^{l^*-1}a^{i-1}\br{n_i^*-n_{i-1}^*}+\br{n-n_{l^*}^*}\cdot\spr{\cCopt}
        \\&\geq
        a^{l^*-1}\br{n-n_{l^*-1}^*}+ \sum_{i=1}^{l^*-1}a^{i-1}\br{\br{n-n_{i-1}^*}-\br{n-n_{i}^*}}+\br{n-n_{l^*}^*}\cdot a^{l^*-1}
        \\&=
        \br{n-n_0^*}+\sum_{i=1}^{l^*-1}\br{a^{i}-a^{i-1}}\cdot\br{n-n_{i}^*}+\br{n-n_{l^*}^*}\cdot a^{l^*-1}
        \\&\geq
        \br{1-\frac{1}{a}}\cdot\sum_{i=0}^{l^*}a^{i}\cdot \br{n-n_i^*}.
        \end{align*}

        Consider any iteration $i\in \brc{1,\dots,l^*}$ of the inner for loop. Let $R_{i-1}=\cU\setminus\cov\br{\angl{\tau_{l^*}^1,\dots,\tau_{l^*}^{i-1}}}$ be the algorithmic residual at the beginning of iteration $i$. The family $\pref{\cCopt}{a^i}$ is precedence-closed, has size at most $a^i\leq a^{i+1}$, and is therefore feasible for the $i$-th $\ProblemBPCSC$ instance. Since it covers $n_i^*$ elements in total and the previously computed solutions cover $n_{i-1}$ elements, it covers at least $n_i^*-n_{i-1}$ elements in $R_{i-1}$. Thus, $\OPT\br{\cI_i}\geq n_i^*-n_{i-1}$.
        By the $\br{\alpha,\beta}$-approximation guarantee, $n_i-n_{i-1}\geq\frac{n_i^*-n_{i-1}}{\beta}$. Rearranging gives $n-n_i\leq\frac{\beta-1}{\beta}\br{n-n_{i-1}}+\frac{1}{\beta}\br{n-n_i^*}$.
        Thus,
        \begin{align*}
            \frac{\br{a-1}\cdot Q}{\alpha}
            &=
            \sum_{i=0}^{l^*}a^{i+3}\cdot \br{n-n_i}
            \\&\leq
            a^3\cdot n + \frac{1}{\beta}\cdot\sum_{i=1}^{l^*}a^{i+3}\cdot \br{n-n_i^*} + \frac{\beta-1}{\beta}\cdot\sum_{i=1}^{l^*}a^{i+3}\cdot \br{n-n_{i-1}}
            \\&\leq
            \frac{2\cdot a^4}{a-1}\cdot\coverageTime\br{\cCopt,\cU}+\frac{\beta-1}{\beta}\cdot\sum_{i=1}^{l^*}a^{i+3}\cdot \br{n-n_{i-1}}
            \\&=
            \frac{2\cdot a^4}{a-1}\cdot\coverageTime\br{\cCopt,\cU}+\frac{\br{\beta-1}\cdot a}{\beta}\cdot\sum_{i=0}^{l^*-1}a^{i+3}\cdot \br{n-n_{i}}
            \\&\leq
            \frac{2\cdot a^4}{a-1}\cdot\coverageTime\br{\cCopt,\cU}+\frac{\br{\beta-1}\cdot a}{\beta}\cdot\br{a-1}\cdot\frac{Q}{\alpha}.
        \end{align*}
        By rearranging the above inequality we obtain:
        \begin{align*}
            Q&\leq \alpha\cdot\coverageTime\br{\cCopt,\cU}\cdot\frac{2\cdot a^4}{\br{a-1}^2\cdot\br{1-\frac{\br{\beta-1}\cdot a}{\beta}}}
            \\&=
            \alpha\cdot\coverageTime\br{\cCopt,\cU}\cdot\frac{2\cdot\br{\frac{3\beta-1}{3\beta-2}}^4}{\br{\frac{3\beta-1}{3\beta-2}-1}^2\cdot\br{1-\frac{\br{\beta-1}\cdot \br{3\beta-1}}{\beta\cdot\br{3\beta-2}}}}
            \\&=
            \alpha\cdot\coverageTime\br{\cCopt,\cU}\cdot\frac{2\cdot \br{3\beta-1}^2\cdot \br{\frac{3\beta-1}{3\beta-2}}^4}{\br{1-\frac{\br{\beta-1}\cdot \br{3\beta-1}}{\beta\cdot\br{3\beta-2}}}}
        \end{align*}
        by our choice of $a$. One can easily check that for $\beta\geq 1$, $\frac{\br{\beta-1}\cdot \br{3\beta-1}}{\beta\cdot\br{3\beta-2}}\leq \frac{\br{\beta-1/3}}{\beta}$ so that:
        \begin{align*}
            Q&\leq
            \alpha\cdot\coverageTime\br{\cCopt,\cU}\cdot\frac{\br{3\beta-1}^2\cdot\br{\frac{2\cdot 3\beta-1}{3\beta-2}}^4}{\br{1-\frac{\br{\beta-1/3}}{\beta}}}
            \\&=
            \alpha\cdot\coverageTime\br{\cCopt,\cU}\cdot 6\beta\cdot \br{3\beta-1}^2\cdot \br{\frac{3\beta-1}{3\beta-2}}^4
            \leq 864\cdot\alpha\cdot\beta^3\cdot\coverageTime\br{\cCopt,\cU}
        \end{align*}
        The lemma follows.
    \end{proof}

    If $\spr{\pi_{l^*}}=m$, then $\pi_{l^*}$ contains every set in $\cS$, and thus covers every element covered by $\cCopt$. The candidate is feasible and its latency is bounded by Lemma~\ref{lem:pi-spread-cost}. Otherwise, let $P_{l^*}$ be the family of sets in $\pi_{l^*}$, let $R_{l^*}=\cU\setminus\cov\br{P_{l^*}}$, and let $r=\spr{\cov\br{P_{l^*},\cU}}$. We then have the following lemmas:
    \begin{lemma}\label{lem:sigma-coverage}
        We have that $\spr{\cov\br{\sigma_{l^*},R_{l^*}}} \geq \frac{\max\brc{0,f\cdot n-r}}{\beta}$ and $\spr{\sigma_{l^*}} = O\br{\alpha}\cdot \spr{\cCopt}$.
    \end{lemma}
    \begin{proof}
        Since $\cCopt$ satisfies $\spr{\cCopt}\leq a^{l^*}$, it is feasible for the final $\ProblemBPCSC$ instance on residual universe $R_{l^*}$. It covers at least $\max\brc{0,f\cdot n-r}$ elements of that residual. The approximation guarantee gives the coverage bound, while $\spr{\sigma_{l^*}}\leq\alpha a^{l^*}\leq\alpha a\spr{\cCopt}=O\br{\alpha}\spr{\cCopt}$ gives the size bound.
    \end{proof}
    Let $\pi_{l^*}^+=\pi_{l^*}\circ\operatorname{ord}\br{\sigma_{l^*}\setminus P_{l^*}}$ denote the sequence obtained by appending only the final solution's new sets.
    \begin{lemma}\label{lem:combined-coverage-cost}
        We have that $\spr{\cov\br{\pi_{l^*}^+,\cU}}\geq \frac{f\cdot n}{\beta}$ and $\coverageTime\br{\pi_{l^*}^+,\cU} = O\br{\alpha\cdot\beta^3}\cdot \coverageTime\br{\cCopt,\cU}$.
    \end{lemma}
    \begin{proof}
        By Lemma~\ref{lem:sigma-coverage}, the prefix and final solution together cover at least
        \[
        r+\frac{\max\brc{0,f\cdot n-r}}{\beta}\geq\frac{f\cdot n}{\beta}
        \]
        elements. Here the residual coverage from $\sigma_{l^*}$ is disjoint from the prefix coverage by definition of $R_{l^*}$.
        For each element $u\in\cU$, let $t_u$ denote the index of the first set in $\pi_{l^*}$ that covers $u$ if $u$ is covered by $\pi_{l^*}$, otherwise set $t_u=\spr{\pi_{l^*}}$.
        Lemma \ref{lem:pi-spread-cost} implies that 
        \[
        \coverageTime\br{\pi_{l^*},\cU}=\sum_{u\in\cU} t_u = O\br{\alpha\cdot\beta^3}\cdot \coverageTime\br{\cCopt,\cU}.
        \]
        Observe, that for each $u\in\cov\br{\pi_{l^*}}$, its cover time in $\pi_{l^*}^+$ is also $t_u$.
        For each $u\notin\cov\br{\pi_{l^*}}$, we have $t_u=\spr{\pi_{l^*}}$. Since this case has $\spr{\pi_{l^*}}<m$, the padding guarantee gives $\spr{\pi_{l^*}}\geq\alpha\beta a^{l^*}$. The appended suffix contains at most $\spr{\sigma_{l^*}}\leq\alpha a^{l^*}\leq\spr{\pi_{l^*}}/\beta$ new sets. Thus the cover time of $u$ in $\pi_{l^*}^+$ is at most $(1+1/\beta)t_u\leq2t_u$.
        As a consequence, due to Lemma~\ref{lem:pi-spread-cost} we get that $\coverageTime\br{\pi_{l^*}^+,\cU} = O\br{\alpha\cdot\beta^3}\cdot \coverageTime\br{\cCopt,\cU}$.
    \end{proof}
By Lemma~\ref{lem:combined-coverage-cost}, the candidate for $l^*$ meets the coverage requirement and has the stated latency bound. Since Algorithm~\ref{alg:fPCMSSC} returns the minimum-latency candidate meeting that coverage requirement, the following theorem follows.
\begin{theorem}\label{thm:BPCSC-to-PCMSSC}
    If there exists a polynomial-time bicriteria $(\alpha,\beta)$-approximation algorithm for $\ProblemBPCSC$, then there exists a polynomial-time bicriteria $\br{\bigo\br{\alpha\cdot\beta^3}, \beta}$-approximation algorithm for $\ProblemPCMSSC$.
\end{theorem}
It should be noted that in all cases when we use Theorem~\ref{thm:BPCSC-to-PCMSSC} we will have $\beta=\bigo\br{1}$, thus yielding a $\br{\bigo\br{\alpha}, \bigo\br{1}}$-approximation algorithm for $\ProblemPCMSSC$.

\subsection{Special cases of precedence constraints} \label{subsection:SpecialCases}
\begin{theorem}\label{thm:BPCSC-inforest}
    There exist polynomial-time bicriteria $\br{1, \frac{e}{e-1}}$-approximation algorithms for $\ProblemBPCSC$ and $\ProblemPCSC$ when the precedence constraints form an inforest.
\end{theorem}
\begin{proof}
    We show how to reduce such an instance to the well-known \textsc{Maximum Coverage} problem (without any precedence constraints) where we allow the non-uniform cost of the sets. In the latter setup including a set $S$ into the cover may impose arbitrary cost $c\br{S}$. For this task a $\br{1, \frac{e}{e-1}}$-approximation algorithm is known \cite{TheBudgetedMaximumCoverageProblem}. Observe that for precedence constraints of the inforest type, for any two sets $S, W\in \cS$ if $\closure{S}\cap\closure{W}\neq \emptyset$, then one is a descendant of the other in the inforest, i.e., either $S\preceq W$ or $W\preceq S$. Thus, we can create an equivalent \textsc{Maximum Coverage} instance in the following way: for each set $S\in \cS$, we create a new set $S'$ in the \textsc{Maximum Coverage} instance, such that $S'=\bigcup_{W\in \closure{S}}W$ and $c\br{S'}=\spr{\closure{S}}$. It is easy to see that any reasonable cover for the \textsc{Maximum Coverage} instance picks at most one representative $S'$ of set $S$ from each tree $T$ in the inforest and this corresponds to picking all of the sets in $\closure{S}$ in the original instance. Thus, this choice corresponds to picking a precedence-closed subtree of this subtree $T$. It is also easy to see, that any precedence-closed subtree of $T$ corresponds to picking a representative $S'$ of some set $S$ in the \textsc{Maximum Coverage} instance. Therefore, any feasible solution to the \textsc{Maximum Coverage} instance corresponds to a feasible solution to $\ProblemBPCSC$ on the original instance with the same cost and coverage, and vice versa. Thus, by Corollary \ref{cor:BPCSC-to-PCSC}, the theorem follows.
\end{proof}

Via a careful examination of Theorem~\ref{thm:BPCSC-to-PCMSSC}, we obtain the following result for $\ProblemPCMSSC$
\begin{corollary}\label{cor:BPCSC-inforest-to-PCMSSC}
    There exists a bicriteria $\br{864\cdot\br{\frac{e}{e-1}}^3+1, \frac{e}{e-1}}\approx\br{3421.7, \frac{e}{e-1}}$-approximation algorithm for $\ProblemPCMSSC$ when the precedence constraints form an inforest running in polynomial time.
\end{corollary}
\begin{theorem}\label{thm:BPCSC-outforest}
    There exist polynomial-time bicriteria $\br{\bigo\br{\log n}, 4}$-approximation algorithms for $\ProblemBPCSC$,  $\ProblemPCSC$ and $\ProblemPCMSSC$ when the precedence constraints form an outforest.
\end{theorem}
\begin{proof}
    We show how to reduce such an instance to $\ProblemGSO$ on tree metrics.
    The latter problem admits a $\br{\bigo\br{\log n}, 4}$-approximation algorithm \cite{ApproxAlgsForOptDTsAndAdapTSPProblems}.
    It should be noted that they give a worse approximation ratio of $\br{\bigo\br{\log^2 n}, 4}$, however their algorithm consists of embedding a general metric into a tree metric with $\bigo\br{\log n}$ distortion.
    Since our metric will be a tree metric, we can skip this step and obtain the improved approximation ratio.
    
    We create a rooted tree as follows. Introduce a dummy root node $r$. For every minimal element $S$ of the outforest, connect $r$ to its representative $v_S$ by an edge of length~$1/2$. For every nonminimal element $W$, let $S$ be its unique immediate predecessor (parent), and connect $v_S$ to $v_W$ by an edge of length~$1/2$. Thus, the path from $r$ to $v_W$ visits all strict predecessors of $W$ before $v_W$. For each element $u\in\cU$, let $X_u=\brc{v_S \mid u\in S}$. For every distinct nonempty set $X$ among the $X_u$, include a GSO group $X$ with weight $w(X)=\spr{\brc{u\in\cU\mid X_u=X}}$.

    Let $\tau$ be any $r$-tour in this tree, and let $\cC_\tau=\brc{S\in\cS\mid v_S\text{ is visited by }\tau}$. The root-to-vertex paths imply that $\cC_\tau$ is precedence-closed. The subtree formed by these paths has exactly $\spr{\cC_\tau}$ edges, one for each visited representative. Since $\tau$ starts and ends at $r$, it traverses each edge in this subtree at least twice. As every edge has length~$1/2$, we have $\spr{\cC_\tau}\leq\spr{\tau}$. The total weight of groups covered by $\tau$ equals $\spr{\cov\br{\cC_\tau,\cU}}$.

    Conversely, for any precedence-closed family $\cC$, the union of the root-to-$v_S$ paths for $S\in\cC$ has exactly $\spr{\cC}$ edges. A tour that traverses each such edge once in each direction visits precisely the representatives in $\cC$, has length $\spr{\cC}$, and covers groups of total weight $\spr{\cov\br{\cC,\cU}}$. Hence, for every budget, tours of at most that length and precedence-closed families of at most that size achieve the same coverage. Thus the reduction preserves the optimum coverage, and by Corollary \ref{cor:BPCSC-to-PCSC} and Theorem \ref{thm:BPCSC-to-PCMSSC} the theorem follows.
\end{proof}


% As of independent interest, we additionally show that if we enforce a certain condition on the input called \emph{$\epsilon$-shallow ancestry}, then for $\epsilon <1$  the greedy solution gives an $\bigo\br{n^\epsilon}$-approximation for $\ProblemMDPCS$.
% \begin{theorem}\label{thm:MDPCS-shallow-ancestry}
% Suppose there exists a constant $C > 0$ and $\epsilon \in (0, 1)$ such that for all $S \in \cS$,
% $
% \spr{\closure{S}} \leq C \cdot
% \spr{\cov\br{\closure{S}}}^\epsilon
% $ ($\epsilon$-shallow ancestry).
% Then a greedy solution provides an $\bigo(n^\epsilon)$-approximation for the $\ProblemMDPCS$ problem.
% \end{theorem}

% \begin{proof}
%  Let $\cC^*$ be an optimal solution and let $\cC$ be a greedy solution (cf. Section~\ref{sec:MDPCS}).
%  Denote for brevity $\delta = \max_{S \in \cS} \Delta\br{\closure{S}}$.
%  There are two cases:
% \begin{enumerate}
%     \item If $\delta \geq n^{1-\epsilon}$, then, we observe that $\Delta\br{\cC} = \delta \geq n^{1-\epsilon}$. Since we can cover at most $n$ elements with at least one set, we have $\Delta\br{\cC^*} \leq n$. Therefore:
% \[
% \frac{\Delta\br{\cC^*}}{\Delta\br{\cC}} \leq \frac{n}{n^{1-\epsilon}} = n^\epsilon
% \]
%     \item Else, if $\delta \leq n^{1-\epsilon}$, we proceed as follows: By definition of density, for any $S \in
%      \cS$, $\spr{\cov\br{\closure{S}}} \leq \delta \cdot \spr{\closure{S}}$. Combining this with the $\epsilon$-shallow ancestry condition, we have that for all $S \in \cS$, $
% \spr{\cov\br{\closure{S}}} \leq \delta \cdot C \cdot \spr{\cov\br{\closure{S}}}^\epsilon
% $. Rearranging this inequality, we get $
% \spr{\cov\br{\closure{S}}} \leq (\delta \cdot C)^{\frac{1}{1-\epsilon}}
% $. We have that:
% \[
% \spr{\cov\br{\cC^*}} = \spr{\bigcup_{S\in\cC^*}\cov\br{\closure{S}}} \leq \sum_{S\in\cC^*} \spr{\cov\br{\closure{S}}} \leq \spr{\cC^*} \cdot (\delta \cdot C)^{\frac{1}{1-\epsilon}}.
% \]
% Therefore,
% \[
% \Delta\br{\cC^*}= \frac{\spr{\cov\br{\cC^*}}}{\spr{\cC^*}} \leq (\delta \cdot C)^{\frac{1}{1-\epsilon}}.
% \]

% By the greedy choice, $\Delta\br{\cC} \geq \delta$ and by assumption $\delta \leq n^{1-\epsilon}$. Thus:
% \[
% \frac{\Delta\br{\cC^*}}{\Delta\br{\cC}} \leq \frac{(\delta \cdot C)^{\frac{1}{1-\epsilon}}}{\delta} = C^{\frac{1}{1-\epsilon}} \cdot \delta^{\frac{\epsilon}{1-\epsilon}} \leq C^{\frac{1}{1-\epsilon}} \cdot (n^{1-\epsilon})^{\frac{\epsilon}{1-\epsilon}} = C^{\frac{1}{1-\epsilon}} \cdot n^\epsilon.
% \]
% \end{enumerate}
% % Since $C$ is constant, theorem follows.
% \end{proof}

% % \subsection{Other Applications}

% % We note that our results for $\ProblemPCSC$ can be applied to the following problem: The \textit{Precedence Constrained Minimum Knapsack Problem} ($\ProblemPCMK$) is as follows: given a set of items $[n]$, each with a cost $c_i$ and a profit $p_i$, a partial order~$\preceq$ on the items, and a value $K$, find a precedence-closed subset of elements $S \subseteq [n]$ such that $p\br{S}=\sum_{i \in S} p_i \geq K$ minimizing $c\br{S}=\sum_{i \in S} c_i$. For this problem we have the following result.
% % \begin{theorem}\label{cor:PCMK}
% %     If there exists a polynomial-time $\br{\alpha, \beta}$-approximation algorithm for $\ProblemPCSC$, then there exists a polynomial-time approximation algorithm for $\ProblemPCMK$ with uniform costs with approximation ratio ${\alpha, \beta}$.
% % \end{theorem}
% % \begin{proof}
% %     Observe that our analysis for $\ProblemPCSC$ assumes that the profit of each element is $1$, however the same argumentation can be applied for the case of non-uniform profits as well. Given a Knapsack instance, we create $n$ sets, each consists of one element with profit equal to the profit of the corresponding item and cost equal to $1$. The precedence constraints are the same as in the Knapsack instance. We run the $\ProblemPCSC$ algorithm on this instance with $f=K/\sum_i p_i$. The claim follows.
% % \end{proof}
% % Observe, that given a Knapsack instance with non-uniform costs we can replace each item $i$ with a chain of $c_i$ items of cost $1$ and profit $0$ for all but the last item with profit $p_i$ and the same precedence constraints. We can then run the $\ProblemPCSC$ algorithm on this instance with $f=K/\sum_i p_i$. Moreover, if we have that $\sum_i p_i\geq n^2$, then one can round the values of profits in the following way to obtain an instance with $O\br{\sum_i p_i\geq n^2}$. Let $T=\frac{\max_{i}p_i}{n}$. Then $p_i'=\floor{\frac{p_i}{T}}$. It is a standard argument that such a rounding increases to optimal value at most $2$ times.
% % We therefore obtain the following corollary:
% % \begin{corollary}
% %     There exist the following polynomial-time approximation algorithms for $\ProblemPCMK$:
% %     \begin{itemize}
% %         \item $\br{1, \frac{e}{e-1}}$-approximation when the precedence constraints form an inforest.
% %         \item $\br{\bigo\br{\log n}, 4}$-approximation when the precedence constraints form an outforest.
% %         \item $\br{\bigo\br{\sqrt{n}\log n}}$-approximation for general precedence constraints.
% %     \end{itemize}
% % \end{corollary}






